\chapter{ODPORÚČANIA PRE PRAX}

Fyzioterapeutický plán u reumatikov musí byť koncipovaný prísne individuálne so zohľadnením aktuálnej aktivity zápalového procesu. Odporúčame systematickú kontrolu klinických indexov, čím sa dosiahne optimálna kalibrácia intenzity interferenčných prúdov a nadväzne kinezioterapie.

Presadzujeme implementáciu včasnej rehabilitácie hneď po stanovení diagnózy. Preventívny charakter včasnej fyzioterapie je zásadný pre potlačenie rozvoja ankylóz a udržanie funkčnej integrity postihnutých kĺbových segmentov.

Považujeme za nevyhnutné edukovať pacientov o význame prístrojovej fyzikálnej terapie. Správna edukácia o protiopuchových a protiestéznych účinkoch IFP môže podstatne zvýšiť aktívnu účasť pacienta na rehabilitačnom procese.

Vďaka preukázateľnému tlmeniu algickej symptomatiky by sa aplikácia interferenčných prúdov mala stať základným pre-kinezioterapeutickým pilierom v dennej rehabilitačnej praxi na reumatologických pracoviskách.

Kombinovaná liečebná schéma (fyzikálna + pohybová) by mala byť štandardizovaná ako súčasť komplexnej starostlivosti pre všetkých pacientov s prejavmi RA na periférnych kĺboch, čím sa akceleruje funkčná obnova motoriky.

Pre dosiahnutie stálosti liečebného progresu odporúčame zavedenie dlhodobých preventívnych cvičení v domácom prostredí. Pacient musí byť motivovaný k pravidelnej autoterapii, ktorá predstavuje kľúčový nástroj k prevencii relapsov a zachovaniu mobility.
